\section{Three inputs to the Hida construction}

% BEGIN F71A1
\paragraph{F71A1. The definite quaternion algebra with its finite adelic splitting.}
For a totally real $F$ of even degree prescribe the local Brauer
invariant $1/2$ at each real place and zero at each finite place.
Their sum is zero in $\mathbb Q/\mathbb Z$. The classification
of quaternion algebras by ramification (Albert--Brauer--Hasse--Noether)
gives a quaternion algebra $D/F$ with exactly those invariants.
It is a division algebra since it is ramified at a real place.
It is central of dimension four and each real completion is
Hamilton's division algebra.

Fix a maximal order. At each finite place choose a splitting of
$D_v$ that carries this order to $M_2(\mathcal O_{F_v})$, using
conjugacy of maximal orders in the split local algebra. These
splittings give one homomorphism $D\to M_2(\mathbb A_{F,f})$.
Its scalar extension over the finite adeles is bijective: a fixed
$F$-basis identifies the restricted product of the $D_v$ with
$D\otimes_F\mathbb A_{F,f}$ because the chosen global order is
the standard basis lattice away from finitely many places. The
local splittings and their inverses preserve the restricted
products. At a real place use the given embedding of $F$ and a
splitting $D_v\otimes_{\mathbb R}\mathbb C\simeq M_2(\mathbb C)$.
The resulting homomorphism on $D$ has precisely the stipulated
action of $F$ on scalars. These are all the fields, including
$\mathrm{jf\_bij}$ and $\mathrm{jc\_alg}$, of
$\mathrm{DefQuatExists}$.
% END F71A1

% BEGIN F71A2
\paragraph{F71A2. Weak unitary descent with the identity lattice at every inert place.}
Take the regular algebraic conjugate self-dual cusp $\Pi$ on
$\mathrm{GL}_n(\mathbb A_E)$ of $\mathrm{DescWeak}$, $2\le n\le4$,
and write $F=E^+$. The extension $E/F$ is unramified at every
finite place. All places over the exceptional set $S$ are split,
so $\Pi$ is spherical at every inert place. The source's unitary
group is that of the identity Hermitian form. It is compact at
infinity and is a pure inner form of the quasi-split unitary
group: Hermitian spaces are obtained by twisting the standard
space by unitary cocycles. At an inert finite place, the group
of the identity lattice is a reductive integral model. It splits
to $\mathrm{GL}_n$ after the unramified quadratic integral
extension, including in residue characteristic two. Its integral
points are the prescribed hyperspecial subgroup.

We use simple generic unitary base change and the multiplicity
formula. More precisely, Labesse, \emph{Changement de base CM et
s\'eries discr\`etes}, Thm.~5.4, gives weak descent of a
cohomological conjugate self-dual cusp. Regular algebraicity
here gives cohomology by the archimedean classification for
cuspidal $\mathrm{GL}_n$ (Clozel). Labesse's degree condition
is automatic: the quadratic Hecke character of $E/F$ is
unramified at every finite place and has the sign character at
each real place. Evaluating on the principal element $-1$ shows
$[F:\mathbb Q]$ is even. In particular we have the standard
base-change parameter and a member on this inner form.
Its global parameter is the simple generic parameter $\Pi[1]$,
by the cuspidality of its transfer and strong multiplicity one
for the isobaric decomposition. The distinction between weak
descent and a prescription at \emph{all} inert places matters
here.

For that prescription use Kaletha--M\'inguez--Shin--White,
\emph{Endoscopic classification of representations: inner forms
of unitary groups}, Thm.~1.7.1 and the scope specified immediately
after it (arXiv:1409.3731). The simple generic parameter and pure
inner twist lie in the case proved there. Its global component
group modulo the invariant center is trivial: a single
conjugate self-dual cuspidal summand of multiplicity one has
centralizer $\{\pm1\}$, already the invariant center.
Thus every restricted product of local packet members satisfying
the usual almost-everywhere normalization occurs. Keep the
archimedean choices. At an inert finite place restriction to
the unramified extension leaves inertia unchanged. The local
parameter downstairs is unramified, since the parameter of the
given spherical $\Pi_w$ has trivial inertia and zero monodromy.
Take the member spherical for the \emph{specified} identity-lattice
hyperspecial subgroup by the unramified Satake correspondence
(M\'inguez, \emph{Unramified representations of unitary groups},
\S\S2--4). These choices agree with the distinguished unramified
choice outside a finite set, as the lattice comes from a single
integral model. At split places use the packet member identified
with the local $\mathrm{GL}_n$ factor under the chosen projection.

The resulting discrete automorphic representation has a copy
in functions on the original adelic group. Take its algebraic
space of smooth finite-adelic and infinity-finite vectors there.
Pure tensors give the original local intertwiners by fixing
nonzero vectors in the other factors. They also give, at each
inert place, a nonzero vector fixed by the indicated subgroup.
Compactness at infinity and smoothness at finite places give
the required continuity, finite level and infinity-finiteness
of every such vector. The restricted tensor product is
irreducible and stable under right translations, so it satisfies
$\mathrm{IsAutRepU}$. At cofinitely many split places the two
specified Satake multisets agree. This proves $\mathrm{DescWeak}$
with all its local conditions.
% END F71A2

% BEGIN F71A3
\paragraph{F71A3. The cohomology entry of the rank-general literature package.}
For the source $\mathrm{LitCoh}$, the finite-module assertions
are G71H1--3, equivalently the finite comparisons recorded in
LD2--3 and DX1--3. For a finite extension
$K\subset\overline{\mathbb Q}_\ell$ the same local duality and
Euler formula apply directly to $G_K$. The character
$\mathrm{cycK}$ is its action on $\mu_{\ell^r}$ and the
exponent is $[K:\mathbb Q_\ell]$. This supplies the three
local-field versions with their original coefficient modules.

It remains to compare $\mathrm{BlochKatoLocal}$ on every supplied
period model. Its $\mathrm{LocTied}$ data identify the two rings
with Fontaine's $B_{\rm cris}$ and $B_{\rm dR}$, intertwine the
Galois actions and identify each filtration step. The concrete
rings in this identification are constructed from
$A_{\rm inf}=W(\mathcal O_{\mathbb C_\ell}^{\flat})$, using the
completed divided powers of the kernel of $\theta$ for
$A_{\rm cris}$ and the $\ker\theta$-adic completion after
inverting $\ell$ for $B_{\rm dR}^+$. Primitive compatible systems
of roots of unity differ by an exponent in $\mathbb Z_\ell^\times$,
so their logarithmic periods are unit multiples of one Fontaine
period $t$. Localizing at all these periods therefore gives
$B_{\rm cris}^+[1/t]$. The fraction field of the de Rham
valuation ring is $B_{\rm dR}$. The maps agree on
$\mathbb Z_\ell$, hence on $\mathbb Q_\ell$ after inverting
$\ell$. Tensoring the supplied ring equivalences with $V$ thus
carries the defining equations for a finite cocycle
$c(g)\otimes1=g x-x$ to the usual crystalline ones and identifies
the filtered de Rham invariants. In particular the source's
dimension condition for $\mathrm{IsDR}$ is the usual de Rham
condition. The cyclotomic character is the one on the specified
compatible roots of unity by the fields of $\mathrm{LocPrime}$.

Use Bloch--Kato, \emph{L-functions and Tamagawa numbers of motives},
\S3. For de Rham $V$ the fundamental exact sequence for the finite
local condition has terms $H^0(K,V)$, $D_{\rm cris}(V)$,
$D_{\rm cris}(V)\oplus D_{\rm dR}(V)/\mathrm{Fil}^0$, and
$H_f^1(K,V)$. Taking dimensions over $\mathbb Q_\ell$ cancels
the two crystalline terms and gives the exact difference of
dimensions in $\mathrm{DimFormula}$. The rational finite local
conditions on $V$ and $V^*(1)$ are mutual annihilators under
local Tate duality (Bloch--Kato, Prop.~3.8). The given
nondegenerate equivariant form identifies the second space with
that very dual. The source's $\mathrm{CobL}$ is equality to the
coboundary of a continuous one-cochain, so zero in $H^2$ gives
the required two cocycle implications.

For the integral statement let $T,T'$ be the supplied dual
lattices. The defining condition on $T'$ says exactly that it
pairs integrally with $T$, hence $V'/T'$ is the discrete Tate
dual of $T$. Local Tate duality pairs the compact group
$H^1(K,T)$ perfectly with $H^1(K,V'/T')$. Compatibility with
the rational pairing identifies the annihilator of the image of
$H_f^1(K,V')$ in the discrete group with the inverse image of
$H_f^1(K,V)$ in the compact group. Indeed, vanishing against
all scalar multiples in the rational subspace modulo
$\mathbb Z_\ell$ is equivalent to zero rational pairing.
Conversely taking annihilators again returns the image, which
is closed in the discrete group. These are exactly the cocycles
named $\mathrm{HfD}$ (rational image plus coboundaries) and
$\mathrm{latticeFiniteCocycles}$ (inverse image). Any change of
the chosen lift of a quotient value changes the cup by a lattice
value, trivial in $\mathbb Q_\ell/\mathbb Z_\ell(1)$.
Here $\mathrm{CobO}$ uses locally constant primitives in this
discrete quotient. Thus its two equivalences also hold for the
given pair and for the specified cyclotomic lattice. This proves
the full $\mathrm{LitCoh}$ at its native period-model quantifier.
% END F71A3
